% LTeX: language=de

\section{Funktionen}

\subsection{Definition}

\begin{Definition}{Funktion}{}
    Wird eine \mycolorbox[yellow]{\textbf{freie Variable}} auf eine davon \mycolorbox[mylightorange]{\textbf{abhängige Größe}} abgebildet, nennt man diese Zuordnung eine \textbf{Funktion}.
    \begin{itemize}[leftmargin=*]
        \item \mycolorbox[yellow]{bestimmt das Experiment} (z. B. Zeit $t$, Volumen $V$, \dots)
        \item \mycolorbox[mylightorange]{kann durch mathematischen Zusammenhang berechnet werden} (z. B. $V = \dfrac{x}{t}$, $x = v \cdot t$, \dots)
    \end{itemize}
    Die Darstellung einer Funktion wird \textbf{Graph} genannt.
\end{Definition}

Funktion und Relation

\subsection{Funktionstypen}

\subsubsection{Lineare Funktionen}

\subsubsection{Quadratische Funktionen}

\subsubsection{Polynome höheren Grads}

\subsection{Ableitung}

\subsubsection{Mitternachtsformel}

Determinante
Scheitel

\subsubsection{Polynomdivision}

\subsubsection{Kettenregel}

\subsubsection{Produktregel}


\subsection{Funktionen mit Parameter}

\subsubsection{Geraden mit Parameter}

\subsubsection{Parabeln mit Parameter}

\subsection{Lineare Gleichungssysteme}